
Вот бутылка с водкой, так называемый спиртуоз. А рядом вы видите Николая Ивановича Серпухова.

Вот из бутылки поднимаются спиртуозные пары. Поглядите, как дышит носом Николай Иванович Серпухов. Видно, ему это очень приятно, и главным образом потому что спиртуоз.

Но обратите внимание на то, что за спиной Николая Ивановича нет ничего. Не то чтобы там не стоял шкап или комод, а совсем ничего нет, даже воздуха нет. Хотите верьте, хотите не верьте, но за спиной Николая Ивановича нет даже безвоздушного пространства, или, как говорится, мирового эфира. Откровенно говоря, ничего нет.

Этого, конечно, и вообразить себе невозможно.

Но на это нам наплевать, нас интересует только спиртуоз и Николай Иванович Серпухов.

Вот Николай Иванович берет рукой бутылку со спиртуозом и подносит ее к своему носу. Николай Иванович нюхает и двигает ртом, как кролик.

Теперь пришло время сказать, что не только за спиной Николая Ивановича, но впереди, так сказать перед грудью и вообще кругом, нет ничего. Полное отсутствие всякого существования, или, как острили когда-то: отсутствие всякого присутствия.

Однако давайте интересоваться только спиртуозом и Николаем Ивановичем.

Представьте себе, Николай Иванович заглядывает внутрь бутылки со спиртуозом, потом подносит ее к губам, запрокидывает бутылку донышком вверх и выпивает, представьте себе, весь спиртуоз.

Вот ловко! Николай Иванович выпил спиртуоз и похлопал глазами. Вот ловко! Как это он!

А мы теперь должны сказать вот что: собственно говоря, не только за спиной Николая Ивановича или спереди и вокруг только, а также и внутри Николая Ивановича ничего не было, ничего не существовало.

Оно, конечно, могло быть так, как мы только что сказали, а сам Николай Иванович мог при этом восхитительно существовать. Это, конечно, верно. Но, откровенно говоря, вся штука в том, что Николай Иванович не существовал и не существует. Вот в чем штука-то.

Вы спросите: <<\textit{А как же бутылка со спиртуозом? Особенно, куда вот делся спиртуоз, если его выпил несуществующий Николай Иванович? Бутылка, скажем, осталась, а где же спиртуоз? Только что был, а вдруг его и нет. Ведь Николай Иванович не существует, говорите вы. Вот как же это так?}>>

Тут мы и сами теряемся в догадках.

А впрочем, что же это мы говорим? Ведь мы сказали, что как внутри, так и снаружи Николая Ивановича ничего не существует. А раз ни внутри, ни снаружи ничего не существует, то значит, и бутылки не существует. Так ведь?

Но с другой стороны, обратите внимание на следующее: если мы говорим, что ничего не существует ни изнутри, ни снаружи, то является вопрос: изнутри и снаружи чего? Что-то, видно, все же существует? А может, и не существует. Тогда для чего же мы говорим изнутри и снаружи?

Нет, тут явно тупик. И мы сами не знаем, что сказать.

До свидания.

\begin{center}
\textit{ВСЕ.}
\end{center}

\begin{flushright}
18 сентября 1934~г.
\end{flushright}

\begin{center}
* * *
\end{center}